\documentclass{article}
\usepackage[utf8]{inputenc}
\usepackage[T1]{fontenc}
\usepackage{geometry}
\geometry{
    a4paper,
    left=2.5cm,
    right=2.5cm,
    top=2.5cm,
    bottom=2.5cm
}
\usepackage{graphicx}
\usepackage{float}
\usepackage{hyperref}
\usepackage[english]{babel}
\usepackage{amsmath}
\usepackage{amssymb}
\usepackage{array}
\usepackage{booktabs}
\usepackage{caption}
\usepackage{subcaption}
\usepackage{enumitem}
\usepackage{fancyhdr}
\usepackage{titlesec}
\usepackage{url}
\usepackage{textcomp}
\usepackage{xcolor}
\usepackage{listings}
\usepackage{multirow}
\usepackage{longtable}
\usepackage{tabularx}
\usepackage{times}

% Set up page style
\pagestyle{fancy}
\fancyhf{}
\rfoot{\thepage}
\renewcommand{\headrulewidth}{0pt}

% Title formatting
\titleformat{\chapter}[display]{\normalfont\bfseries}{\chaptertitle\thechapter}{0pt}{\Huge}
\titlespacing*{\chapter}{0pt}{50pt}{40pt}

% Document info
\title{Template Document}
\author{Academic Support Agent}
\date{\today}

\begin{document}
\maketitle

% % % % % % % % % % % % % % % % % % % % % % % % % % % % % % % % % % % % % % 
\pretextual  %%%% GERA AS PAGINAS PRE-TEXTUAIS   
% % % % % % % % % % % % % % % % % % % % % % % % % % % % % % % % % % % % % % 

% % % % % % % % % % % % % % % % % % % % % % % % % % % % % % % % % % % % % % 
% % % % % CORPO DO TRABALHO - INCLUA OS SEUS TEXTOS AQUI
% % % % % SUGESTAO -> UTILIZE ARQUIVOS EXTERNOS A PARTIR DO COMANDO \input




% % % % % % % % % % % % % % % % % % % % % % % % % % % % % % % % % % % % % % 
% % % % % % % % % % INICIO DAS PAGINAS TEXTUAIS % % % % % % % % % % % % % % 
% % % % % % % % % % % % % % % % % % % % % % % % % % % % % % % % % % % % % % 




% % % % % % % % % % % % % % % % % % % % % % % % % % % % % % % % % % % % % % 
% % % % % % % % % % % % % INTRODUCAO % % % % % % % % % % % % % % % % % % % % 
% % % % % % % % % % % % % % % % % % % % % % % % % % % % % % % % % % % % % % 
\introducao{%
\par Insira aqui a introdu o!!!

\par Insira aqui a introdu o!!!

\par Insira aqui a introdu o!!!

\par Insira aqui a introdu o!!!

\par Insira aqui a introdu o!!!

% \par Insira aqui a introdu o!!!

}%
% % % % % % % % % % % % % % % % % % % % % % % % % % % % % % % % % % % % % % 
\geraintro  %%%% GERA INTRODUCAO   % % % % % % % % % % % % % % % % % % % % % 
% % % % % % % % % % % % % % % % % % % % % % % % % % % % % % % % % % % % % % 
% % % % % % % % % % % % % % % % % % % % % % % % % % % % % % % % % % % % % % 

\chapter{Algumas notas sobre bibliografia}


\par O arquivo \textit{referencias.bib}   o nome padr o das refer ncias para este modelo. Para alter -lo modifique o argumento entre chaves na defini o de bibliografia.

   \section{Modelos de cita o}
      \par No arquivo modelo de refer ncias, tamb m existem alguns exemplos de diferentes classes de cita es. Todas elas podem ser usadas com o \textit{$\backslash$cite$\{$label$\}$} ou \textit{$\backslash$citeonline$\{$label$\}$}, dependendo da forma de cita o\footnote{Este   um teste de nota de rodap }.

	\par Al m disso, pode-se incluir obras na bibliografia que nortearam o trabalho mesmo que elas n o apare am diretamente no texto, utilizando o comando \textit{$\backslash$nocite$\{$label$\}$}. Al m disso, pode-se citar v rios trabalhos em conjunto, por exemplo:
	
	\begin{center}\rule{0.5\textwidth}{1pt}\\$\backslash cite\{label1,label2,label3,...\}$\end{center}

	\begin{verbatim}
	Os ventos do norte n o movem moinhos \cite{tcc:mintegui2014, 
	diss:anabor2004, tese:anabor2008, livro:halliday28ed, 
	livro:fedorova:v1, site:amsglo:fog}.
	\end{verbatim}
	    
	\par Os ventos do norte\footnote{Este   um teste de nota de rodap } n o movem moinhos \cite{tcc:mintegui2014,diss:anabor2004,tese:anabor2008,livro:halliday28ed,livro:fedorova:v1,site:amsglo:fog}.
        
        \begin{center}\rule{0.5\textwidth}{1pt}\\$\backslash citeonline\{label1,label2,label3,...\}$\end{center}
        
	\begin{verbatim}
	Segundo \citeonline{cap:livro:djuric1994, livro:aris1989, 
	livro:cotton1989, cap:livro:wyngaard1981, livro:fedorova:v1, 
	artigo:fujita1981,puhalesBLT2010,artigo:janjic2002,
	cbmet:anabor2012,site:wrfhome}, os ventos do norte n o movem moinhos.
	\end{verbatim}
            
        \par Segundo \citeonline{cap:livro:djuric1994, livro:aris1989, livro:cotton1989, cap:livro:wyngaard1981, livro:fedorova:v1, artigo:fujita1981,puhalesBLT2010,artigo:janjic2002,cbmet:anabor2012,site:wrfhome}, os ventos\footnote{Este   um teste de nota de rodap } do norte n o movem moinhos.
        \begin{center}\rule{0.5\textwidth}{1pt}\end{center}   
         \par OBS: Trabalhos com mais de tr s autores aparecer o com a abreviatura ``et al'' no corpo do texto, por m todos os nomes da bibliografia (norma da UFSM que   definida nas op es do abntcite nas defini es do arquivo). 
  
  \subsection{O comando \textit{apud}}
  \par A cita o \textit{apud} ocorre quando voc  cita algum autor atrav s de outra obra, sem ter consultado-a propriamente. Neste caso a cita o   feita da seguinte forma:
  \begin{center}
  \rule{0.5\textwidth}{1pt}\\
  $\backslash apud\{material\_lido\}\{material\_citado\_no\_material\_lido\}$ \\
  \end{center}
\begin{verbatim}
Sobre a circula o geral da atmosfera pode-se dizer que os ventos do norte
n o movem moinhos \apud{livro:monin:v1}{apud:richardson1922}.
\end{verbatim}
  
  Sobre a circula o geral da atmosfera pode-se dizer que os ventos do norte n o movem moinhos \apud{livro:monin:v1}{apud:richardson1922}.
\begin{center}\rule{0.5\textwidth}{1pt}\end{center}  
  \par Nesse caso, na bibliografia s  constar  a obra consultada e n o aquele referenciada pela obra. Para que isso ocorra naturalmente, a obra consultada deve ser inclu da normalmente no arquivo referencias.bib enquanto a obra referenciada indiretamente deve ser inclu da com a op o \textit{@hidden}, conforme o modelo de refer ncias.

      \subsubsection{\textit{Apud on line}}
      
      \par O \textit{apudonline} se aplica da mesma maneira que o \textit{apud} descrito anteriormente. O termo \textit{on line}   alusivo ao \textit{$\backslash$citeonline$\{$label$\}$} definido no abntex. Nesse caso a cita o   feita da seguinte forma:
      \begin{center}
      \rule{0.5\textwidth}{1pt}
            $\backslash apudonline\{material\_lido\}\{material\_citado\_no\_material\_lido\}$ \\
	    \end{center}

 \begin{verbatim}
Segundo \apudonline{livro:monin:v1}{apud:richardson1922}, os ventos do
norte n o movem moinhos.
\end{verbatim}

            Segundo \apudonline{livro:monin:v1}{apud:richardson1922}, os ventos do norte n o movem moinhos.

       \paragraph{Teste de se o quin ria}
       
       
       \par Texto texto texto.
       
       
         \chapter{Quadros, tabelas e figuras}
         
         \par Na MDT da UFSM h  uma clara diferen a entre tabelas e quadros, quanto a sua apresenta o. Aqui, para inserir tabelas usa-se o ambiente tradicionalmente definido \textit{table}. A partir deste modelo simples:
         

        
        
         \begin{verbatim}
		\begin{table}[ht]
		\centering
		\caption{Modelo de tabela para MDT-UFSM.}
		\begin{tabular}{ c c c }
		\hline
		Abacate & Banana & Canela \\
		\hline
		21 & 34 & 56 \\
		-3 & 245 & 23 \\
		-25 & -0,57 & 2 \\
                \hline
                 \end{tabular}
                \vspace{\baselineskip} %%% linha em branco para atender a norma
		\fonte{Adaptado de \citeonline{livro:halliday28ed}.}
		\end{table}
	  \end{verbatim}
         
         \noindent resulta:
         
         \begin{table}[ht]
         \centering
         \caption{Modelo de tabela para MDT-UFSM.}
         \begin{tabular}{ c c c }
         \hline
         Abacate & Banana & Canela \\
         \hline
         21 & 34 & 56 \\
         -3 & 245 & 23 \\
         -25 & -0,57 & 2 \\
         \hline
         \end{tabular}
         \vspace{\baselineskip} %%% linha em branco para atender a norma
          \fonte{Adaptado de \citeonline{livro:halliday28ed}.}
         \end{table}
         
         \par Note que, adicionalmente, foi definido um comando novo: ``fonte''. Ele serve para indicar a fonte da tabela quando necess rio, mas tamb m pode ser usado em outros ambientes.
         
         \par Para inserir quados, foi criado um novo ambiente: ``quadro''. Ele seve ser usado de forma semelhante a tabela, como o ambiente tabular. Contudo, neste caso, as linhas verticais e horizontais est o sempre presentes. Um exemplo simples   o seguinte: 
         
         
         \begin{verbatim}
	    \begin{quadro}
   	    \caption{Modelo de quadro para MTD-UFSM.}
	    \centering
	    \begin{tabular}{| c |c |c }
	    \hline
	    Abacate & Banana & Canela \\
	    \hline
	    21 & 34 & 56 \\
	    \hline
	    -3 & 245 & 23 \\
	    \hline
	    -25 & -0,57 & 2 \\
	    \hline
	    \end{tabular}
	    \vspace{\baselineskip} %%% linha em branco para atender a norma
	    \fonte{Adaptado de \citeonline{livro:halliday28ed}.}
	    \end{quadro}
         \end{verbatim}
         
         \noindent resultando:
         
	    \begin{quadro}
	    \caption{Modelo de quadro para MTD-UFSM.}
	    \centering
	    \begin{tabular}{| c |c |c |}
	    \hline
	    Abacate & Banana & Canela \\
	    \hline
	    21 & 34 & 56 \\
	    \hline
	    -3 & 245 & 23 \\
	    \hline
	    -25 & -0,57 & 2 \\
	    \hline
	    \end{tabular}
	    \vspace{\baselineskip} %%% linha em branco para atender a norma
    	    \fonte{Adaptado de \citeonline{livro:halliday28ed}.}
	    \end{quadro}
	    

         \noindent Assim como para as tabelas e figuras, j  est  definida uma lista de quadros. Al m disso, o comando ``fonte'' tamb m pode ser usado aqui se necess rio, bem como para figuras. Vale lembrar que, na MDT-UFSM, as legendas de quadros e figuras aparecem embaixo das mesmas. J  nas tabelas, em cima. A fonte sempre embaixo.
         
         \par As figuras devem ser inseridas com o ambiente padr o: \textit{figure}. Veja um exemplo simples:
         
         \begin{verbatim}
	    \begin{figure}[ht]
		\centering
		\includegraphics[width=0.6\textwidth]{figuras/pretextuais.png}
		\caption{\label{exepretex} Sequ ncia dos elementos pr -textuais da MDT-UFSM}
		\vspace{\baselineskip} %%% linha em branco para atender a norma
		\fonte{Adaptado de \citeonline{man:MDTUFSM2012}.}
	    \end{figure}
         \end{verbatim}
         
         \begin{figure}[ht]
     	    \caption{\label{exepretex} Sequ ncia dos elementos pr -textuais da MDT-UFSM}
	    \centering
	    \includegraphics[width=0.6\textwidth]{figuras/pretextuais.png}
	    \vspace{\baselineskip} %%% linha em branco para atender a norma
            \fonte{Adaptado de \citeonline{man:MDTUFSM2012}.}
         \end{figure}
        
        
         
\chapter{Conclus o}

	\par Conclus o do trabalho. 

         \begin{verbatim}
	    \begin{ilustracao}[ht]
		\centering
		\includegraphics[width=0.6\textwidth]{figuras/pretextuais.png}
		\caption{\label{exepretex} Sequ ncia dos elementos pr -textuais da MDT-UFSM}
		\vspace{\baselineskip} %%% linha em branco para atender a norma
		\fonte{Adaptado de \citeonline{man:MDTUFSM2012}.}
	    \end{ilustracao}
         \end{verbatim}
         
         \begin{ilustracao}[ht]
     	    \caption{\label{exepretex} Sequ ncia dos elementos pr -textuais da MDT-UFSM}
	    \centering
	    \includegraphics[width=0.6\textwidth]{figuras/pretextuais.png}
	    \vspace{\baselineskip} %%% linha em branco para atender a norma
            \fonte{Adaptado de \citeonline{man:MDTUFSM2012}.}
         \end{ilustracao}
% % % % % % % % % % % % % % % % % % % % % % % % % % % % % % % % % % % % % % 
% % % % % % % % % % % % FIM DAS PAGINAS TEXTUAIS % % % % % % % % % % % % % % 
% % % % % % % % % % % % % % % % % % % % % % % % % % % % % % % % % % % % % % 



% % % % % % % % % % % % % % % % % % % % % % % % % % % % % % % % % % % % % % 	
% % % % % % % % % % % % % BIBLIOGRAFIA  % % % % % % % % % % % % % % % % % % 
% % % % % % % % % % % % % % % % % % % % % % % % % % % % % % % % % % % % % % 	

\bibliografia{referencias}  %%%%% BIBLIOGRAFIA -> INCLUIR NAS CHAVES O NOME DO ARQUIVO *.BIB	
	
	
	
% % % % % % % % % % % % % % % % % % % % % % % % % % % % % % % % % % % % % 	
% % % % % % % % % % % % % APENDICES % % % % % % % % % % % % % % % % % % %
% % % % % % % % % % % % % % % % % % % % % % % % % % % % % % % % % % % % % 	
	\apendice %%%% TEXTOS A PARIR DESTE PONTO SERAO CONSIDERADOS APENDICES

%\chapter{Demonstra o de algo}
%        \par Algo como ap ndice.  

        
        
        
% % % % % % % % % % % % % % % % % % % % % % % % % % % % % % % % % % % % % % 	
% % % % % % % % % % % % % % % ANEXOS  % % % % % % % % % % % % % % % % % % % 
% % % % % % % % % % % % % % % % % % % % % % % % % % % % % % % % % % % % % % 	
        \anexo    %%%% TEXTOS A PARIR DESTE PONTO SERAO CONSIDERADOS ANEXOS
        
%\chapter{Algo interessante que algu m fez}
%         \par Algo como anexo.
%         
%          \begin{verbatim}
%	    \begin{ilustracao}[ht]
%		\centering
%		\includegraphics[width=0.6\textwidth]{figuras/pretextuais.png}
%		\caption{\label{exepretex} Sequ ncia dos elementos pr -textuais da MDT-UFSM}
%		\vspace{\baselineskip} %%% linha em branco para atender a norma
%		\fonte{Adaptado de \citeonline{man:MDTUFSM2012}.}
%	    \end{ilustracao}
%         \end{verbatim}
         
%         \begin{grafico}[ht]
%     	    \caption{\label{exepretex1} Sequ ncia dos elementos pr -textuais da MDT-UFSM}
%	    \centering
%	    \includegraphics[width=0.6\textwidth]{figuras/pretextuais.png}
%	    \vspace{\baselineskip} %%% linha em branco para atender a norma
%            \fonte{Adaptado de \citeonline{man:MDTUFSM2012}.}
%         \end{grafico}
         
         
%         \section*{Teste de se o dentro do anexo}
\end{document}